\paragraph{In-context learning.} Since \citet{brown2020language} demonstrated the in-context learning ability of GPT-3, there has been a significant interest in improving and understanding this capability \citep{liu2021makes, min2021noisy, zhao2021calibrate, lu2021fantastically, rubin2021learning, min2021metaicl, chen2021meta, mishra2021reframing, lampinen2022can}.
The works most relevant to ours are as follows.
\citet{xie2022an} propose a Bayesian inference framework explaining how in-context learning works despite formatting differences between training and inference distributions. \citet{razeghi2022impact} show that in-context learning for numerical reasoning tasks is better for instances whose terms are more prevalent in training data. \citet{min2021noisy} demonstrate tasks where in-context learning  works  even when the prompt outputs are chosen randomly, questioning to what extent these models are truly  learning new tasks on-the-fly, while \citet{rong_2021} 
gives examples of novel tasks on which these models demonstrate on-the-fly learning ability. \citet{chan2022data} demonstrate that distributional properties such as long-tailedness are crucial for in-context learning on an image-based few-shot dataset. \citet{olsson2022context} and  \citet{elhage2021mathematical} consider a different framing of in-context learning, referring to any model behavior that utilizes information in a prompt to make predictions that improve with prompt size. They hypothesize the existence of special circuits inside Transformer models responsible for in-context learning, that can complete prompts by copying previous similar patterns in the prompt sequence. 
\citet{pesut2022who} and \citet[Table~16]{dinh2022lift} consider in-context learning for small tabular datasets and learning problems in one and two dimensions, and show that GPT-3 can obtain non-trivial accuracy.
Our work contributes to and complements this line of work, by posing in-context learning as a well-defined problem of learning function classes at inference time, and empirically investigating training models that in-context learn simple function classes.

\paragraph{Transformers.} There is a long line of work investigating the capabilities~\citep{vaswani2017attention,dehghani2018universal,yun2019transformers,perez2019turing,yao2021self,bhattamishra2020computational,zhang2022unveiling}, limitations~\citep{hahn2020theoretical, bhattamishra2020ability}, applications~\citep{lu2021pretrained,dosovitskiy2020image, parmar2018image}, and internal workings~\citep{elhage2021mathematical,snell2021approximating,weiss2021thinking,edelman2022inductive,olsson2022context} of Transformer models.
Most similar to our work, \citet{muller2021transformers} and \citet{nguyen2022transformer} demonstrate the ability of Transformer models to solve prediction tasks using the input context, albeit in different settings. \citet{muller2021transformers} introduce a ``Prior-data fitted transformer network'' that is trained to approximate Bayesian inference with priors such as Gaussian processes and Bayesian neural networks,  and use it to perform downstream tasks such as tabular dataset classification and few-shot image classification. \citet{nguyen2022transformer} introduce Transformer neural processes, building on prior work on neural processes~\citep{garnelo2018neural, garnelo2018conditional, kim2019attentive}, and show that they achieve state-of-the art performance on tasks such as image completion and contextual multi-armed bandits.
Our work complements these works, focusing on understanding the in-context learning ability of Transformers for various simple function classes and the extent to which this ability extrapolates beyond the training distribution.

\paragraph{Meta learning.} Training a model to perform in-context learning
can be viewed as an instance of the more general learning-to-learn or meta-learning paradigm \citep{schmidhuber1987evolutionary,naik1992meta,thrun2012learning}.
Typical approaches from this extensive line of work (see \citep{hospedales2020meta} for a survey) include:
training a meta-learner on how to update the parameters of a downstream learner \citep{bengio1995optimization, li2016learning}, learning parameter initializations from which one can quickly train for many downstream tasks \citep{finn2017model, ravi2016optimization},
learning latent embeddings that allow for effective similarity search~\citep{snell2017prototypical}.
Most relevant to our setting are approaches that directly  take as input examples from a downstream task and a query input and produce the corresponding output~\citep{hochreiter2001learning,mishra2018simple, santoro2016meta, garnelo2018neural, garnelo2018conditional,kirsch2021meta}.
Our work contributes to this line of work, by investigating the learning-to-learn abilities of Transformer models in a well-defined setting.

\paragraph{Data-driven algorithm design.}
Another line of work aims to discover algorithms that perform well on a distribution of inputs~\citep{horvitz2001bayesian,xu2008satzilla,vinyals2015pointer,bello2016neural,khalil2017learning,selsam2018learning,schwarzschild2021can} (as opposed to algorithms with guarantees on their worst-case performance).
See~\citet{balcan2020data} for a survey on advancements on the theoretical foundations of such algorithms.
Our work can be viewed as part of this line of work, as we train Transformer models to discover algorithms for different learning problems.